\documentclass[11pt,a4paper]{article}

% ==============================================================================
% CONFIGURACIÓN DE IDIOMA Y CODIFICACIÓN
% ==============================================================================
\usepackage[utf8]{inputenc}
\usepackage[T1]{fontenc}
\usepackage[spanish,es-tabla]{babel}

% ==============================================================================
% GEOMETRÍA Y DISEÑO DE PÁGINA
% ==============================================================================
\usepackage{geometry}
\geometry{
    top=2.5cm,
    bottom=2.5cm,
    left=2.5cm,
    right=2.5cm,
    headheight=22pt,
    footskip=1.5cm
}

% ==============================================================================
% PALETA DE COLORES CORPORATIVOS
% ==============================================================================
\usepackage[table]{xcolor}
\definecolor{PrimaryNavy}{RGB}{15, 44, 89}        % Azul marino profundo corporativo
\definecolor{SecondarySteel}{RGB}{30, 96, 145}     % Azul acero / acento profesional
\definecolor{AccentAmber}{RGB}{217, 119, 6}        % Ámbar para alertas y notas
\definecolor{BgSoftGray}{RGB}{248, 250, 252}       % Fondo gris claro para cajas
\definecolor{BorderGray}{RGB}{203, 213, 225}       % Borde sutil
\definecolor{DarkText}{RGB}{30, 41, 59}            % Texto principal carbón
\definecolor{MutedText}{RGB}{100, 116, 139}        % Texto secundario / metadatos

% ==============================================================================
% TIPOGRAFÍA Y MICRO-TIPOGRAFÍA
% ==============================================================================
\usepackage{lmodern}
\renewcommand{\familydefault}{\sfdefault} % Fuente Sans-Serif moderna ejecutiva
\usepackage{microtype}

% ==============================================================================
% PAQUETES TABULARES Y LISTAS
% ==============================================================================
\usepackage{booktabs}
\usepackage{tabularx}
\usepackage{multirow}
\usepackage{enumitem}
\setlist{nosep, leftmargin=1.8em, labelsep=0.5em}

% ==============================================================================
% CAJAS RESALTADAS (TCOLORBOX)
% ==============================================================================
\usepackage[many]{tcolorbox}

% Estilo para Historias de Usuario
\newtcolorbox{UserStoryBox}[2]{%
    enhanced,
    colback=BgSoftGray,
    colframe=PrimaryNavy,
    coltitle=white,
    fonttitle=\bfseries\sffamily,
    title={#1 \hfill \small\textbf{ID: #2}},
    arc=3mm,
    boxrule=1pt,
    titlerule=0pt,
    top=8pt,
    bottom=8pt,
    left=12pt,
    right=12pt,
    drop shadow=black!8!white,
    before skip=12pt,
    after skip=12pt
}

% Estilo para Integraciones
\newtcolorbox{IntegrationBox}[2]{%
    enhanced,
    colback=white,
    colframe=SecondarySteel,
    coltitle=white,
    fonttitle=\bfseries\sffamily,
    title={#1 \hfill \small\textbf{Ref: #2}},
    arc=2mm,
    boxrule=1.2pt,
    titlerule=0pt,
    top=8pt,
    bottom=8pt,
    left=12pt,
    right=12pt,
    shadow={1.5mm}{-1.5mm}{0mm}{black!6!white},
    before skip=12pt,
    after skip=12pt
}

% Estilo para Glosario / Notas Técnicas
\newtcolorbox{GlossaryBox}[1]{%
    enhanced,
    colback=BgSoftGray,
    colframe=BorderGray,
    coltitle=PrimaryNavy,
    fonttitle=\bfseries\sffamily,
    title={#1},
    attach boxed title to top left={yshift=-2mm, xshift=5mm},
    boxed title style={colback=white, colframe=SecondarySteel, boxrule=0.8pt, arc=1.5mm},
    arc=2mm,
    boxrule=0.8pt,
    top=12pt,
    bottom=8pt,
    left=12pt,
    right=12pt,
    before skip=10pt,
    after skip=10pt
}

% ==============================================================================
% ENCABEZADOS Y PIES DE PÁGINA (FANCYHDR)
% ==============================================================================
\usepackage{fancyhdr}
\usepackage{lastpage}

\pagestyle{fancy}
\fancyhf{}
\renewcommand{\headrulewidth}{0.6pt}
\renewcommand{\footrulewidth}{0.4pt}
\renewcommand{\headrule}{\hbox to\headwidth{\color{SecondarySteel}\leaders\hrule height \headrulewidth\hfill}}
\renewcommand{\footrule}{\hbox to\headwidth{\color{BorderGray}\leaders\hrule height \footrulewidth\hfill}}

\fancyhead[L]{\small\sffamily\textcolor{PrimaryNavy}{\textbf{CRM NextGen}} $\cdot$ Captura de Requerimientos}
\fancyhead[R]{\small\sffamily\textcolor{MutedText}{Versión 1.0}}
\fancyfoot[L]{\small\sffamily\textcolor{MutedText}{Documento Confidencial de Ingeniería de Software}}
\fancyfoot[R]{\small\sffamily\textcolor{PrimaryNavy}{\textbf{Página \thepage\ de \pageref{LastPage}}}}

% Redefinir estilo plain para páginas de inicio de capítulo o índice
\fancypagestyle{plain}{
    \fancyhf{}
    \fancyfoot[L]{\small\sffamily\textcolor{MutedText}{Documento Confidencial de Ingeniería de Software}}
    \fancyfoot[R]{\small\sffamily\textcolor{PrimaryNavy}{\textbf{Página \thepage\ de \pageref{LastPage}}}}
    \renewcommand{\headrulewidth}{0pt}
}

% ==============================================================================
% HIPERVÍNCULOS E ÍNDICE INTERACTIVO
% ==============================================================================
\usepackage{hyperref}
\hypersetup{
    colorlinks=true,
    linkcolor=PrimaryNavy,
    urlcolor=SecondarySteel,
    citecolor=SecondarySteel,
    filecolor=SecondarySteel,
    plainpages=false,
    pdfpagelabels=true,
    pdftitle={Documento de Captura de Requerimientos - Sistema CRM},
    pdfauthor={Equipo de Arquitectura e Ingeniería de Software},
    pdfsubject={Especificación de Requerimientos y Diseño Funcional},
    pdfkeywords={CRM, Requerimientos, Historias de Usuario, Integraciones, RNF},
    bookmarksnumbered=true,
    pdfstartview={FitH}
}

% ==============================================================================
% PERSONALIZACIÓN DE TÍTULOS Y SECCIONES
% ==============================================================================
\usepackage{titlesec}
\titleformat{\section}
    {\color{PrimaryNavy}\normalfont\Large\bfseries\sffamily}
    {\thesection.}{0.8em}{}
    [\vspace{0.2em}\color{SecondarySteel}\titlerule]

\titleformat{\subsection}
    {\color{SecondarySteel}\normalfont\large\bfseries\sffamily}
    {\thesubsection}{0.6em}{}

% Espaciado de párrafos
\setlength{\parindent}{0pt}
\setlength{\parskip}{6pt}

% ==============================================================================
% CUERPO DEL DOCUMENTO
% ==============================================================================
\begin{document}

% ------------------------------------------------------------------------------
% PORTADA EJECUTIVA
% ------------------------------------------------------------------------------
\hypersetup{pageanchor=false}
\begin{titlepage}
    \pagecolor{white}
    \begin{tikzpicture}[remember picture, overlay]
        % Barra vertical decorativa izquierda
        \fill[color=PrimaryNavy] (current page.north west) rectangle ++(1.2cm, -\paperheight);
        \fill[color=SecondarySteel] ([xshift=1.2cm]current page.north west) rectangle ++(0.3cm, -\paperheight);
    \end{tikzpicture}
    
    \vspace*{1.5cm}
    \begin{flushleft}
        \hspace*{1.2cm}
        \begin{minipage}{0.88\textwidth}
            {\large\sffamily\textcolor{SecondarySteel}{\textbf{ESPECIFICACIÓN DE SOFTWARE Y SISTEMAS EMPRESARIALES}}}\\[0.4cm]
            {\Huge\sffamily\bfseries\textcolor{PrimaryNavy}{Documento de Captura de}}\\[0.2cm]
            {\Huge\sffamily\bfseries\textcolor{PrimaryNavy}{Requerimientos}}\\[0.4cm]
            {\LARGE\sffamily\bfseries\textcolor{SecondarySteel}{Sistema CRM Corporativo}}\\[0.8cm]
            
            \noindent\textcolor{SecondarySteel}{\rule{\linewidth}{1.5pt}}\\[0.8cm]
            
            {\large\sffamily\textcolor{DarkText}{\textbf{Subtítulo:}} Historias de Usuario, Integraciones y Requerimientos No Funcionales}\\[1.5cm]
            
            % Ficha Técnica de Portada
            \begin{tcolorbox}[
                colback=BgSoftGray,
                colframe=BorderGray,
                arc=2mm,
                boxrule=1pt,
                left=15pt,
                right=15pt,
                top=12pt,
                bottom=12pt
            ]
                \sffamily
                \begin{tabular}{ll}
                    \textbf{\textcolor{PrimaryNavy}{Versión:}} & 1.0 \\[4pt]
                    \textbf{\textcolor{PrimaryNavy}{Fecha:}} & \today \\[4pt]
                    \textbf{\textcolor{PrimaryNavy}{Estado:}} & Aprobado para Fase de Desarrollo \\[4pt]
                    \textbf{\textcolor{PrimaryNavy}{Área Responsable:}} & Dirección de Tecnología e Innovación \\[4pt]
                    \textbf{\textcolor{PrimaryNavy}{Líder de Producto:}} & Equipo de Ingeniería de Software
                \end{tabular}
            \end{tcolorbox}
            
            \vfill
            
            {\small\sffamily\textcolor{MutedText}{
                Este documento define la línea base de especificación funcional y técnica del Sistema CRM.
                Su distribución está sujeta a las directivas de confidencialidad institucional.
            }}
        \end{minipage}
    \end{flushleft}
\end{titlepage}
\hypersetup{pageanchor=true}

\newpage
\pagenumbering{roman}

% ------------------------------------------------------------------------------
% TABLA DE CONTENIDOS
% ------------------------------------------------------------------------------
\vspace*{0.5cm}
{
    \hypersetup{linkcolor=PrimaryNavy}
    \tableofcontents
}

\vfill

\begin{tcolorbox}[
    enhanced,
    colback=BgSoftGray,
    colframe=BorderGray,
    boxrule=0.8pt,
    arc=2mm,
    left=12pt,
    right=12pt,
    top=8pt,
    bottom=8pt
]
    \sffamily\small
    \textcolor{PrimaryNavy}{\textbf{Estructura y Trazabilidad de Requerimientos:}}\\
    El presente documento está estructurado en cinco secciones clave que cubren desde el control de versiones y necesidades operativas de los usuarios, hasta los criterios de aceptación en integraciones críticas y requerimientos no funcionales de arquitectura.
\end{tcolorbox}

\vspace{0.4cm}
\begin{center}
    \small\sffamily\textcolor{MutedText}{Documento de Requerimientos de Software $\cdot$ Sistema CRM Corporativo}
\end{center}

\newpage
\pagenumbering{arabic}

% ------------------------------------------------------------------------------
% SECCIÓN 1: HISTORIAL DE VERSIONES
% ------------------------------------------------------------------------------
\section{Historial de Versiones}

A continuación se detalla el control de cambios y revisiones formales aplicadas al presente documento de requerimientos:

\vspace{0.3cm}
\begin{table}[h]
    \centering
    \sffamily\small
    \renewcommand{\arraystretch}{1.3}
    \begin{tabularx}{\textwidth}{c c p{4.2cm} X}
        \toprule
        \rowcolor{PrimaryNavy}
        \textcolor{white}{\textbf{Versión}} & 
        \textcolor{white}{\textbf{Fecha}} & 
        \textcolor{white}{\textbf{Autor}} & 
        \textcolor{white}{\textbf{Descripción del Cambio}} \\
        \midrule
        \rowcolor{BgSoftGray}
        \textbf{0.1} & 10/10/2026 & Equipo de Análisis Funcional & Elaboración inicial del levantamiento de historias de usuario. \\
        \textbf{0.5} & 15/10/2026 & Arquitectura de Integraciones & Incorporación de especificaciones de API RESTful y OAuth 2.0. \\
        \rowcolor{BgSoftGray}
        \textbf{0.9} & 18/10/2026 & Especialista en Seguridad TI & Definición de RNF (RBAC, GDPR, TLS 1.3 y latencias de servicio). \\
        \textbf{1.0} & \today & Comité de Aprobación Técnica & Consolidación final y liberación para arquitectura y desarrollo. \\
        \bottomrule
    \end{tabularx}
    \caption{Registro oficial de versiones y trazabilidad de cambios.}
    \label{tab:historial_versiones}
\end{table}

\vspace{0.6cm}

% ------------------------------------------------------------------------------
% SECCIÓN 2: HISTORIAS DE USUARIO
% ------------------------------------------------------------------------------
\section{Historias de Usuario (User Stories)}

Las Historias de Usuario definen las necesidades y funcionalidades del sistema desde la perspectiva de los diferentes perfiles operativos del CRM, aplicando el estándar estructurado \textit{Rol -- Acción -- Beneficio}.

% HU-001
\begin{UserStoryBox}{Gestión Centralizada de Contactos}{HU-001}
    \sffamily
    \begin{itemize}[label=\textcolor{SecondarySteel}{\textbullet}]
        \item \textbf{Como:} Ejecutivo de Ventas.
        \item \textbf{Quiero:} Registrar, consultar y actualizar la información de clientes (datos de contacto, historial de interacciones y estado de la cuenta).
        \item \textbf{Para:} Mantener un seguimiento continuo y personalizado de cada prospecto.
    \end{itemize}
    \tcblower
    \small\textcolor{MutedText}{\textbf{Prioridad:} Alta \hfill \textbf{Iteración sugerida:} Sprint 1}
\end{UserStoryBox}

% HU-002
\begin{UserStoryBox}{Gestión del Flujo de Ventas (Pipeline)}{HU-002}
    \sffamily
    \begin{itemize}[label=\textcolor{SecondarySteel}{\textbullet}]
        \item \textbf{Como:} Director Comercial.
        \item \textbf{Quiero:} Visualizar y gestionar las oportunidades de negocio en un tablero por etapas (prospección, negociación, cierre).
        \item \textbf{Para:} Proyectar los ingresos y dar seguimiento al rendimiento del equipo.
    \end{itemize}
    \tcblower
    \small\textcolor{MutedText}{\textbf{Prioridad:} Crítica \hfill \textbf{Iteración sugerida:} Sprint 1}
\end{UserStoryBox}

% HU-003
\begin{UserStoryBox}{Automatización de Tareas y Recordatorios}{HU-003}
    \sffamily
    \begin{itemize}[label=\textcolor{SecondarySteel}{\textbullet}]
        \item \textbf{Como:} Asesor de Servicio al Cliente.
        \item \textbf{Quiero:} Programar tareas, llamadas y seguimientos automáticos vinculados a cada cliente.
        \item \textbf{Para:} Evitar omisiones en la atención y mejorar el tiempo de respuesta.
    \end{itemize}
    \tcblower
    \small\textcolor{MutedText}{\textbf{Prioridad:} Media-Alta \hfill \textbf{Iteración sugerida:} Sprint 2}
\end{UserStoryBox}

% ------------------------------------------------------------------------------
% SECCIÓN 3: INTEGRACIONES Y CRITERIOS DE ACEPTACIÓN
% ------------------------------------------------------------------------------
\section{Integraciones y Criterios de Aceptación}

El sistema CRM interactuará con el ecosistema tecnológico empresarial para asegurar continuidad operativa y sincronización de datos sin duplicidad.

% INT-01
\begin{IntegrationBox}{Integración con Correo Electrónico (Gmail / Outlook)}{INT-01}
    \sffamily
    \textbf{\textcolor{PrimaryNavy}{Descripción del Alcance:}}\\
    Sincronización bidireccional de correos electrónicos para registrar la comunicación directa dentro de la plataforma CRM, manteniendo trazabilidad completa de cada intercambio con los clientes.
    
    \vspace{0.2cm}
    \textbf{\textcolor{PrimaryNavy}{Criterios de Aceptación:}}
    \begin{enumerate}[label=\textbf{\arabic*.}]
        \item \textbf{Autenticación Segura:} Conexión delegada mediante protocolo \textbf{OAuth 2.0} sin almacenamiento de credenciales directas en la base de datos del CRM.
        \item \textbf{Asociación Automática:} Vinculación instantánea de correos enviados y recibidos al historial unificado del contacto con base en la dirección de e-mail.
        \item \textbf{Gestión de Archivos Adjuntos:} Habilitación de envío y almacenamiento de documentos adjuntos directamente desde la interfaz del CRM a través del cliente de correo enlazado.
    \end{enumerate}
\end{IntegrationBox}

% INT-02
\begin{IntegrationBox}{Integración con ERP / Sistema de Facturación}{INT-02}
    \sffamily
    \textbf{\textcolor{PrimaryNavy}{Descripción del Alcance:}}\\
    Conexión mediante arquitectura \textbf{API RESTful} para sincronizar en tiempo real o en lotes el estado de facturación, emisión de pedidos y conciliación de pagos.
    
    \vspace{0.2cm}
    \textbf{\textcolor{PrimaryNavy}{Criterios de Aceptación:}}
    \begin{enumerate}[label=\textbf{\arabic*.}]
        \item \textbf{Generación de Factura Borrador:} Creación automatizada del borrador de factura en el ERP inmediatamente al marcar una oportunidad como \textit{«Venta Cerrada»} en el CRM.
        \item \textbf{Sincronización de Pagos:} Actualización bidireccional del estado financiero y balance de pagos directamente visible en la ficha del cliente en el CRM.
        \item \textbf{Resiliencia y Monitoreo:} Implementación de cola de reintentos exponenciales (\textit{exponential backoff}) y emisión de alertas críticas al administrador ante anomalías o fallos de conexión con el endpoint del ERP.
    \end{enumerate}
\end{IntegrationBox}

\newpage

% ------------------------------------------------------------------------------
% SECCIÓN 4: REQUERIMIENTOS NO FUNCIONALES (RNF)
% ------------------------------------------------------------------------------
\section{Requerimientos No Funcionales (RNF)}

Los Requerimientos No Funcionales estipulan las condiciones de calidad, seguridad, eficiencia operativa y experiencia de usuario que el sistema debe satisfacer de forma estricta.

\vspace{0.3cm}
\begin{table}[h]
    \centering
    \sffamily\small
    \renewcommand{\arraystretch}{1.4}
    \begin{tabularx}{\textwidth}{c >{\raggedright\arraybackslash}p{2.8cm} >{\raggedright\arraybackslash}X >{\raggedright\arraybackslash}p{3.8cm}}
        \toprule
        \rowcolor{PrimaryNavy}
        \textcolor{white}{\textbf{ID}} & 
        \textcolor{white}{\textbf{Categoría}} & 
        \textcolor{white}{\textbf{Especificación Técnica}} & 
        \textcolor{white}{\textbf{Métrica / Criterio}} \\
        \midrule
        \rowcolor{BgSoftGray}
        \textbf{RNF-01} & 
        \textbf{Rendimiento} & 
        Tiempos de respuesta óptimos en renderizado de vistas y procesamiento en capa de servicios API. & 
        Carga de vistas $< 2\,\text{s}$. \newline Respuestas API $< 500\,\text{ms}$. \\
        
        \textbf{RNF-02} & 
        \textbf{Seguridad y Cumplimiento} & 
        Cifrado integral de canales en tránsito, control granular de privilegios y observancia de normativas de protección de datos. & 
        HTTPS con TLS 1.3, modelo RBAC y cumplimiento estricto GDPR. \\
        
        \rowcolor{BgSoftGray}
        \textbf{RNF-03} & 
        \textbf{Disponibilidad y Escalabilidad} & 
        Garantía de continuidad del servicio en entornos productivos y capacidad elástica ante picos de demanda. & 
        Uptime del $99.5\%$ y soporte para $+50\%$ de usuarios concurrentes. \\
        
        \textbf{RNF-04} & 
        \textbf{Usabilidad y Accesibilidad} & 
        Diseño ergonómico y adaptativo (\textit{responsive}) optimizado para la agilidad en la gestión diaria de prospectos. & 
        Multi-dispositivo. Alta de prospecto en menos de $3$ clics. \\
        \bottomrule
    \end{tabularx}
    \caption{Matriz de Requerimientos No Funcionales del Sistema CRM.}
    \label{tab:rnf}
\end{table}

\vspace{0.5cm}

% ------------------------------------------------------------------------------
% SECCIÓN 5: GLOSARIO Y NOTAS TÉCNICAS
% ------------------------------------------------------------------------------
\section{Glosario y Notas Técnicas}

Para garantizar un entendimiento inequívoco entre los equipos de desarrollo, analistas de negocio y los tomadores de decisión, se definen los siguientes términos técnicos clave:

\begin{GlossaryBox}{CRM (\textit{Customer Relationship Management})}
    Plataforma de software orientada a la gestión integrada de las relaciones comerciales, ventas, seguimiento de prospectos y atención al cliente, cuyo fin es consolidar toda la interacción en una base unificada.
\end{GlossaryBox}

\begin{GlossaryBox}{ERP (\textit{Enterprise Resource Planning})}
    Sistema de planificación de recursos empresariales encargado de centralizar y automatizar procesos contables, facturación, compras, inventarios y logística dentro de una organización.
\end{GlossaryBox}

\begin{GlossaryBox}{API RESTful (\textit{Representational State Transfer})}
    Interfaz de comunicación basada en los estándares del protocolo HTTP (GET, POST, PUT, DELETE), con arquitectura sin estado (\textit{stateless}) y payloads usualmente formateados en JSON, diseñada para interoperabilidad entre sistemas distribuidos.
\end{GlossaryBox}

\begin{GlossaryBox}{OAuth 2.0}
    Protocolo de autorización estándar de la industria que permite a aplicaciones de terceros obtener acceso limitado y seguro a servicios de usuario (por ejemplo, Google Workspace o Microsoft 365) mediante \textit{tokens}, prescindiendo del almacenamiento de contraseñas.
\end{GlossaryBox}

\begin{GlossaryBox}{RBAC (\textit{Role-Based Access Control})}
    Mecanismo de seguridad y control de acceso basado en roles que restringe o habilita operaciones en el sistema de acuerdo con las responsabilidades y privilegios asignados al perfil del usuario.
\end{GlossaryBox}

\begin{GlossaryBox}{Uptime}
    Métrica porcentual que mide el tiempo de disponibilidad efectiva e ininterrumpida de una infraestructura de software o servicio en la nube durante un periodo determinado (e.g., $99.5\%$ anual o mensual).
\end{GlossaryBox}

\vspace{0.8cm}
\begin{center}
    \small\sffamily\textcolor{MutedText}{--- Fin del Documento de Especificación de Requerimientos ---}
\end{center}

\end{document}
